\clearpage
\section{Realisierungsphase}
\label{sec:Realisierungsphase}
Zur Dokumentierung dieser Phase wird eine Art "`Bottom-up-Prinzip"' genutzt. Dazu werden technisch oder fachlich getrennte
Teilbereiche näher beleuchtet und dann in Zusammenhang gebracht.



\subsection{Projektstruktur}
\label{sub:Projektstruktur}
\footnotetext{Internet Explorer (bis Version 7 pflicht), Mozilla Firefox, Google Chrome, Opera.}
\begin{wrapfigure}[14]{r}[25pt]{0\textwidth}
\tikzstyle{every node}=[draw=black,thick,anchor=west]
\tikzstyle{new}=[dashed,fill=hig]
\scalebox{0.8}{
\begin{tikzpicture}[%
  grow via three points={one child at (0.5,-0.7) and
  two children at (0.5,-0.7) and (0.5,-1.4)},
  edge from parent path={(\tikzparentnode.south) |- (\tikzchildnode.west)}]
\texttt{
\node {com.<mandant>.contrax}
    child { node {party.domain}
    	child { node [new] {repository}
    	}
    }	
    child [missing] {}	
    child { node {shared.infrastructure.async}
    	child { node [new] {atmosphere}
    		child { node [new] {domain}
				child { node [new] {model}
				}
				child { node [new] {service}
					child { node [new] {impl}
					}
				}
			}
		}
	};
}
\end{tikzpicture}
}
\caption{Differenz zur vorherigen Packetstruktur} 
\label{app:DifferenzPacketstruktur} 
\end{wrapfigure}
Es wird die im VMS verwendete Paketstruktur übernommen. Diese ist angelehnt an das "`Domain Driven
Design"'\footnote{Definition siehe \citep{DDD}.} aufgebaut. Die Struktur startet unterhalb von \Code{"`com.<mandant>.contrax"'},
dabei steht \Code{<mandant>} für das Kürzel eines Kunden. Auf der obersten Ebene befinden sich die fachlichen Säulen, gefolgt von
den technischen Layern. Eine Erklärung dieser ist in einer tabellarischen und auszugsweisen Form im \AnhangK{app:Paketstruktur} zu
finden. Durch die Implementierung des Benachrichtigungssystems wurden die in \Abbildung{app:DifferenzPacketstruktur} gestrichelten
Packete hinzugefügt. Der Großteil der neuen Dateien konnte aber in die bisherige Struktur eingegliedert werden.



\subsection{Einrichten/Konfigurieren von Atmosphere}
\label{sub:Einrichten/Konfigurieren von Atmosphere}
Das Framework wird wie im \Anhang{app:pom.xml} gezeigt per Maven\footnote{Maven ist ein "`Build-Management-Tool"' basierend auf
Java.} in der Version 2.2.6 eingebunden. In \Abkuerzung{VMS} werden alle Servlet-Container Komponenten (wie Filter und Listener)
in der \Klasse{VmsWebAppInitializer} konfiguriert. Hier wird die code-basierte Alternative, einer Servlet 3.0+ Umgebung, zu
einer klassischen \Code{web.xml}-Konfiguration genutzt. Nachdem das Atmosphere-Servlet im Webapplikationskontext unter dem Pfad
\Code{"/notify/*"} registriert wurde, wird es über eine \Fachbegriff{Spring}-annotierte Konfigurationsklasse (siehe
\Anhang{app:AtmosphereConfig}) entsprechend eingerichtet. Für eine asynchrone Verarbeitung muss das Atmosphere-Servlet sowie alle
unter diesem Pfad wirksamen Filter und Listener mit dem Attribut "`asyncSupported=true"' konfiguriert werden. (siehe
\Anhang{app:VmsWebAppInitializer}).

\subsection{Kommunikation: Server zu Client}
\label{sub:Kommunikation: Server zu Client}
In diesem Abschnitt wird hauptsächlich der Code vom \Anhang{app:NotificationServiceHandler} erläutert und sich
auf \Anhang{app:main.js} bezogen. Außerdem werden unter Anderem einige Java-Klassen aus dem Framework Atmosphere namentlich
genannt und daher im Folgenden erklärt. Für einen besseren Lesefluss wird bei den Begriffen der Präfix "`Atmosphere"' im weiteren
Verlauf vernachlässigt.
\begin{description}
	\item [AtmosphereHandler] entspricht einem Servlet. Jeder \Klasse{AtmosphereHandler} ist einem konkretem Pfad unterhalb von
	\Code{"/notify/*"} zugeordnet und übernimmt im Allgemeinen die Handhabung von POST/GET-Requests. 
	\item [AtmosphereResource] ist eine Java(-Wrapper)-Klasse die einen HTTP-Request, also ein Verbindung, darstellt. Über die
	jeweilige Instanz kann eine Verbindung \zB angehalten und Informationen über diese abgerufen werden.
	\item [AtmosphereSession] ermöglicht im Backend die persistente Speicherung von Informationen über eine Verbindung. Jede
	\Klasse{AtmosphereResource} kann eine \Klasse{AtmosphereSession} haben, welche auch nur solange lebt, wie die
	\Klasse{AtmosphereResource} selbst.
	\item [Broadcaster] sind Nachrichtenverteiler. Über die letztlich "`Events"' mittels der Methode \Methode{broadcast} übermittelt
	werden können. Jede \Klasse{AtmosphereResource} Instanz wird einem \Klasse{Broadcaster} zugeteilt, welcher dann \zB über
	\MethodenRef{AtmosphereResource\#getBroadcaster()} erreichbar ist. Zusätzlich kann eine \Klasse{AtmosphereResource} mehrere
	\Klasse{Broadcaster} abonnieren. Nur die abonnierten \Klasse{Broadcaster} wissen dann von der \Klasse{AtmosphereResource}, nicht
	andersrum. Über eine entsprechende \sog \Code{Factory}\footnote{Eine Factory kapselt die Logik zur Erstellung eines Objektes von
	der eigentlichen Business-Logik ab. An die \Code{Factory}-Methode übergebene Daten werden dann zur Instanziierung eines
	bestimmten Objekttyps genutzt, meist ist dies eine abgeleitete Klasse \zB eines Interfaces (zu Deutsch: Schnittstelle).
	\citep[vgl.][]{Factory}} können \Klasse{Broadcaster} ermittelt werden.
\end{description}

\subsubsection{Verbindungshaltung}
\label{subsub:Verbindungshaltung}
Bei jedem Seitenwechsel stellt der Client per JavaScript eine zusätzliche Anfrage (\Code{\acs{XML}HttpRequest}) über den
\Code{"/notify/*"} Pfad an den Server. Diese \Fachbegriff{Resource} wird vom entsprechenden \Fachbegriff{Handler} jedoch nicht
beantwortet, sondern bis auf weiteres "`suspended"'. Die Verbindung wird daher aufrecht erhalten und ermöglicht dem Server im
späteren Verlauf den Client/Browser bei Bedarf zu benachrichtigen. Wird ein Event veröffentlicht, werden die entsprechenden
\Fachbegriff{Resource}s beantwortet. Unmittelbar nach Erhalt dieser Antwort, sendet der Client erneut eine Anfrage. Dieses
Vorgehen wurde bereits als "`long-polling"' im \VerweisVL{sub:Soll-Konzept} aufgezeigt.\\
\Fachbegriff{Atmosphere} stellt bei der Verwaltung dieser \Fachbegriff{Resource}s sicher, dass es nur eine \Fachbegriff{Resource}
pro Browser gibt. Unter Anderem durch die Vergabe einer einmaligen Id für jeden Browser.

\subsubsection{Identifikation einer Verbindung}
\label{subsub:Identifikation einer Verbindung}
Nachdem eine \Fachbegriff{Resource} suspendiert wurde, beschafft die \Klasse{AtmosphereResourceSessionFactory} die zugehörige
\Fachbegriff{Session}. Mittels \Fachbegriff{Spring} wird der Benutzer ermittelt und seine Id in der \Fachbegriff{Session}
gespeichert. Im späteren Verlauf kann so eine \Fachbegriff{Resource} dem jeweiligen Benutzer zugeordnet werden. Dies ist nötig,
wenn \zB explizit nur ein bestimmter Benutzer benachrichtigt werden soll.

Soll eine Benutzergruppe über ein Event informiert werden, so wird zur Identifikation der Verbindungen eine Reihe von 
\Fachbegriff{Broadcaster}n genutzt. Bei der Initialisierung vom \Klasse{NotificationServiceHandler} wird dazu eine
Map\footnote{Liste von Schlüssel/Wert-Paaren.} von \Fachbegriff{Broadcaster}n pro Benutzergruppe instanziiert. Anschließend wird
eine Liste von allen dem System bekannten Benutzergruppen erstellt. Zuletzt wird für jeden Eintrag ein \Fachbegriff{Broadcaster},
mit dem Gruppennamen als Kennung, definiert und der Map \Code{userRoleBroadcasterMap} hinzugefügt. Nun kann der
Nachrichtenverteiler für eine Benutzergruppe mittels der \Klasse{BroadcasterFactory} und der Kennung ermittelt werden.

\subsection{Datenhaltung}
\label{sub:Datenhaltung}
An dieser Stelle wird aufgezeigt, wie die Webanwendung bestimmte Informationen abspeichert. Es wird auf viele Anhänge aus dem
\AnhangKapitel{app:Quellcode} und zu Teilen aus dem \AnhangKapitel{app:Weboberfläche} Bezug genommen.

\subsubsection{UserStatusRepository}
\label{subsub:UserStatusRepository}
\begin{description}
	\item [Repository]
Ein Repository (zu Deutsch: Ablage, Aufbewahrungsort) stellt ein Mittel dar, die Geschäftslogik von der Beschaffung
und/oder der Instanziierung zu trennen. Erzeugte Objekte liegen dann im Repository abrufbereit. Häufig genutzte Lese-, Schreib-
und Änderungsvorgänge werden oftmals direkt vom Repository zur Verfügung gestellt. \citep[vgl.][]{Repository}
\end{description}

Zum Zeitpunkt der Implementierung wird bisher nur der letzte Besuch der Offenen-Task-Übersicht eines jeden Benutzers im
\Fachbegriff{Repository} gespeichert. Die Ermittlung des Zeitstempels wird unten stehend erklärt. Da es sich lediglich um ein
\Code{DateTime}-Objekt handelt, könnte diese Information in der vorhandenen Datenbank gespeichert werden. Gemäß dem Abschnitt
Datenhaltung aus \VerweisVL{sub:Technisches Konzept} wurde jedoch das \Fachbegriff{Repository} entwickelt.

\paragraph{Letzter Besuch der Offenen-Task-Übersicht}
\label{para:Offene-Task-Übersicht}
Beim Aufruf dieser Übersicht wird zunächst der aktuelle "`zuletzt besucht"'-Zeitstempel beim \Klasse{UserStatusRepository}
erfragt. Existiert für den genannten Benutzer kein Zeitstempel, so wird der Zeitpunkt "`\Code{01.01.1970
00:00:00}"'\footnote{Datum-Bedeutung: \url{http://openbook.rheinwerk-verlag.de/javainsel9/javainsel_15_002.htm}
(Aufruf:~26.03.15).} angelegt. Der Zeitstempel wird dann für die Markierung neuer Tasks in der Weboberfläche genutzt.\\
Da ein Aufruf der genannten Seite gleichzeitig diesen Zeitstempel verändern soll, musste dem JSF-Lebenszyklus\footnote{Erklärung:
\url{http://jsfatwork.irian.at/book_de/jsf.html\#!idx:/jsf.html:2.6} (Aufruf:~26.03.15).} ein "`PostRenderView"'-Event hinzugefügt
werden. Dieses Event wird erst ausgelöst, wenn eine Seite fertig gebaut ist und somit der bisherige Zeitstempel von der
Weboberfläche verwendet wurde. Wurde das Event auf der besagten Seite (\Code{listOpenTasks.xhtml}) ausgelöst wird die Methode
\MethodenRef{OpenTaskListController\#visitedOpenTasks()} aufgerufen, die den Prozess zur Aktualisierung des Zeitstempels anstößt.

\subsubsection{TaskNotificationStatisticsModel}
\label{subsub:TaskNotificationStatisticsModel}
Diese Java-Klasse hält die Daten bereit, die für die Darstellung der Status-Anzeige\VerweisVF{sub:Soll-Konzept} benötigt werden
und wird von der Weboberfläche\footnote{In diesem Kontext von \Anhang{app:masterTemplate}.} über den
\Klasse{NotificationController} genutzt. Dieses "`Model"' ist ein Ergebnis der \Abkuerzung{SQL}-Abfrage in
\Methode{getTaskNotificationModel(VmsUser user, DateTime lastVisit)} der Klasse \Klasse{TaskQueryService}. Dabei werden
Informationen über die Task-Typen Fehler und Benutzerinteraktion abgefragt. Dazu gehört die Anzahl offener Tasks eines Typs und
ob jeweils neue Tasks seit dem letzten Besuch der Offenen-Tasks-Übersicht hinzugekommen sind. Jede Information wird in einer
Spalte angelegt. Anhand der Spaltenbezeichnung wird die Information entsprechend in das
\Klasse{TaskNotificationStatisticsModel}-Objekt zugeordnet. Dabei handelt es sich entweder um einen numerischen oder booleschen
Wert.


\subsection{Datenübertragung}
\label{sub:Datenübertragung}
Bei derzeitiger Implementierung, nach \VerweisVL{sub:Technisches Konzept}, muss jede Event-Klasse die Klasse
\Klasse{AtmosphereEvent} erweitern, und kann \ggfs die Methode \Methode{toJSON()} überschreiben.\\
Eine Beispielimplementierung stellt dabei die Klasse \Klasse{TaskNotificationEvent} dar. In dieser wird ersichtlich wie Daten
pro Task-Typ abgelegt werden können. Es wurde darauf verzichtet, \\die \Methode{AtmosphereEvent\#toJSON()} anzupassen. Dieses
Event informiert den Client lediglich über neue Tasks, enthält aber nicht die zur Anzeige benötigten Informationen (Anzahl der offenen
Task pro Typ). Diese werden wie in Kapitel \VerweisVL{sub:Event-Handling} beschrieben anschließend vom Client abgefragt.
Die direkte Übertragung der Informationen hat für Probleme gesorgt. Zur Einhaltung des Zeitrahmens wurde der o.g. Lösungsansatz
gewählt.


\subsection{Event-Handling}
\label{sub:Event-Handling}
In diesem Kapitel wird aufgezeigt, wie Events im Backend publiziert werden können und wie diese im Frontend behandelt werden. Es
wird Bezug auf \Anhang{app:TaskNotificationService}, \Anhang{app:main.js}, \Anhang{app:NotificationController} und
\Anhang{app:main.less} genommen.

\subsubsection{Im Backend}
\label{subsub:Backend}
Im Backend wird ein Service (\Klasse{TaskNotificationService}) bereitgestellt, welcher zur Publizierung von
Task-Benachrichtigungs-Events über \Fachbegriff{Atmosphere} genutzt wird.\\
Ein konkreter Service zum Versenden von Events implementiert die im Interface \Klasse{NotificationService} deklarierte Methode
\Methode{pushNotification}. Diese Methode ist mehrfach überladen. Je nach Ziel kann der Methode als Parameter einen Benutzer
oder eine Benutzergruppe übergeben werden. Bei Ersterem reicht die Benutzer-Id. Es gibt eine \Klasse{AbstractNotificationService}
Klasse mit wichtigen Hilfsmethoden die von jedem derartigen Service benötigt werden.

\subsubsection{Im Frontend}
\label{subsub:Frontend}
Beim Client wird das Event-Handling über den \Methode{onMessage}-Listener geregelt. Siehe \Anhang{app:main.js}. Alle Events, die
vom Server veröffentlicht werden, kommen bei erfolgreicher Übermittlung bei dieser Methode an. Zunächst wird versucht, den
Nachrichteninhalt in ein gültiges \Abkuerzung{JSON}-Objekt umzuwandeln. Falls die Konvertierung scheitert, wird bei bisheriger
Implementierung die Nachricht vom Server nicht neu versendet. War die Umwandlung erfolgreich, so wird zur Unterscheidung der
Event-Typen die Variable \Code{event} mit einem entsprechenden String verglichen. Dazu wird Kenntnis über die serverseitige
Enum-Definition dieser Strings benötigt.\\
Liegt ein Task-Benachrichtigungs-Event vor, wird mittels \Fachbegriff{jQuery} ein Klick auf einem versteckten Button in der
Weboberfläche ausgelöst. Dadurch wird der Prozess zur \Abkuerzung{UI}-Aktualisierung angestoßen. Genauer wird zunächst das
dahinterliegende Model der Status-Anzeigen zurückgesetzt, anschließend wird ein "`partial Ajax-Request"' ausgelöst. Dieser fordert
das genannte Model beim Server an, da es zurückgesetzt wurde, wird es von diesem neu gebaut. Zum Schluss wird das nun
aktualisierte Model zum Rendern der neuen Status-Anzeige genutzt.

\subsection{UI-Entwurfsumsetzung}
\label{sub:UI-Entwurfsumsetzung}
Die spezifischen und konkreten Ideen aus \VerweisKL{sub:Soll-Konzept} und \VerweisKL{sub:UI-Entwurf} wurden umgesetzt. Der
Task-Typ Benutzerinteraktion erscheint in einem blauen Oval, während die Status-Anzeige der Error-Tasks im roten Rechteck
angezeigt wird. Durch die Formen allein sind die Task-Typen bereits unterscheidbar. Zusätzlich besitzt jede Status-Anzeige einen
\sog Tooltip\footnote{Wenn die Maus über eine Element gehalten wird, erscheint eine kleine Infobox neben der Maus.}, welcher kurz
und schlüssig die Anzeigen differenziert und erklärt. Dadurch kann auf eine Anwenderdokumentation verzichtet werden.

In der Task-Übersicht ist erkennbar, dass außerdem alle weiteren Task-Typen einem Pastellton zugeordnet wurden. Darin ist
ebenfalls der Hover-Effekt\footnote{Wenn die Maus über eine Element gehalten wird, wird dieses Element gesondert hervorgehoben.}
sichtbar, der zur Verschönerung der Status-Anzeigen und des Farbbalkens hinzugefügt wurde.
Angelehnt am bereits vorhandenem Hover-Effekt der Tabelle, bei der eine Zeile, über der die Maus liegt, subtil farblich
hervorgehoben wird. So wird auch der farbliche Balken und die Status-Anzeige in der Farbtiefe angehoben (Siehe farblicher Balken
in \Anhang{app:UI-Entwurf}).

